\section{Bảng tổng hợp công thức và ký hiệu}

Bảng dưới đây tóm tắt các ký hiệu và công thức quan trọng được sử dụng xuyên suốt báo cáo. Mục tiêu của phần này là giúp người đọc tra cứu nhanh các biểu thức chính trước khi đi vào phần chứng minh và phân tích thuật toán.


\subsection{Bảng tổng hợp ký hiệu}

\renewcommand{\arraystretch}{1.22}
\setlength{\tabcolsep}{6pt}
\begin{table}[H]
\centering
\small
\caption{Tổng hợp các ký hiệu chính trong báo cáo}
\begin{tabular}{|>{\raggedright\arraybackslash}p{3.2cm}|>{\raggedright\arraybackslash}p{4.2cm}|>{\raggedright\arraybackslash}p{6.6cm}|}
\hline
\textbf{Nhóm} & \textbf{Ký hiệu} & \textbf{Ý nghĩa} \\
\hline
Dữ liệu & $m$ & Số lượng điểm dữ liệu huấn luyện. \\
\hline
Dữ liệu & $N$ & Số chiều đặc trưng của dữ liệu đầu vào. \\
\hline
Dữ liệu & $\bm{x}_i$ & Điểm dữ liệu thứ $i$. \\
\hline
Dữ liệu & $y_i$ & Nhãn đúng của điểm dữ liệu thứ $i$ trong bài toán đơn nhãn. \\
\hline
Dữ liệu & $\mathcal{X}$ & Không gian đầu vào. \\
\hline
Dữ liệu & $\mathcal{Y}$ & Tập nhãn/lớp của bài toán. \\
\hline
Dữ liệu & $k=|\mathcal{Y}|$ & Số lượng lớp trong bài toán phân loại đa lớp. \\
\hline
Dữ liệu & $L$ & Số nhãn trong bài toán extreme multi-label classification. \\
\hline
Hàm điểm và lề & $h(\bm{x},y)$ & Hàm điểm số của mẫu $\bm{x}$ đối với nhãn $y$. \\
\hline
Hàm điểm và lề & $g_h(\bm{x})$ & Hàm dự đoán nhãn cuối cùng sinh bởi hàm điểm $h$. \\
\hline
\end{tabular}
\end{table}

\begin{table}[H]
\centering
\small
\caption{Tổng hợp các ký hiệu chính trong báo cáo (tiếp)}
\begin{tabular}{|>{\raggedright\arraybackslash}p{3.2cm}|>{\raggedright\arraybackslash}p{4.2cm}|>{\raggedright\arraybackslash}p{6.6cm}|}
\hline
\textbf{Nhóm} & \textbf{Ký hiệu} & \textbf{Ý nghĩa} \\
\hline
Hàm điểm và lề & $\rho_h(\bm{x},y)$ & Lề đa lớp tại mẫu $(\bm{x},y)$. \\
\hline
Hàm điểm và lề & $\rho$ & Ngưỡng lề mong muốn trong hàm mất mát lề. \\
\hline
Hàm điểm và lề & $\Phi_\rho$ & Hàm mất mát lề với tham số lề $\rho$. \\
\hline
Hàm điểm và lề & $R(h)$ & Sai số thật/kỳ vọng của giả thuyết $h$. \\
\hline
Hàm điểm và lề & $\widehat{R}_\rho(h)$ & Sai số lề kinh nghiệm trên tập huấn luyện. \\
\hline
Rademacher & $S$ & Tập mẫu dùng để tính độ phức tạp Rademacher kinh nghiệm. \\
\hline
Rademacher & $\sigma_i$ & Biến Rademacher độc lập, nhận giá trị $\{-1,+1\}$. \\
\hline
Rademacher & $\widehat{\mathfrak{R}}_S(\mathcal{F})$ & Độ phức tạp Rademacher kinh nghiệm của lớp hàm $\mathcal{F}$. \\
\hline
Rademacher & $\mathfrak{R}_m(\mathcal{F})$ & Độ phức tạp Rademacher kỳ vọng theo mẫu kích thước $m$. \\
\hline
Rademacher & $\Pi_1(\mathcal{H})$ & Lớp hàm tọa độ thu được khi cố định một nhãn của $h\in\mathcal{H}$. \\
\hline
\end{tabular}
\end{table}
\begin{table}[H]
\centering
\small
\caption{Tổng hợp các ký hiệu chính trong báo cáo (tiếp)}
\begin{tabular}{|>{\raggedright\arraybackslash}p{3.2cm}|>{\raggedright\arraybackslash}p{4.2cm}|>{\raggedright\arraybackslash}p{6.6cm}|}
\hline
\textbf{Nhóm} & \textbf{Ký hiệu} & \textbf{Ý nghĩa} \\
\hline
Kernel/SVM & $\Phi(\bm{x}_i)$ & Ánh xạ đặc trưng của $\bm{x}_i$ vào không gian Hilbert. \\
\hline
Kernel/SVM & $K(\bm{x}_i,\bm{x}_j)$ & Hàm kernel, bằng tích vô hướng của hai ánh xạ đặc trưng. \\
\hline
Kernel/SVM & $\bm{w}_l$ & Vector trọng số tương ứng với lớp $l$. \\
\hline
Kernel/SVM & $W$ & Ma trận trọng số ghép từ $(\bm{w}_1,\ldots,\bm{w}_k)$. \\
\hline
Kernel/SVM & $\xi_i$ & Biến lơi lỏng biểu diễn mức vi phạm lề tại mẫu $\bm{x}_i$. \\
\hline
Kernel/SVM & $C$ & Siêu tham số điều chuẩn trong Multi-class SVM. \\
\hline
Kernel/SVM & $\Lambda$ & Chặn trên chuẩn của ma trận trọng số $W$. \\
\hline
Kernel/SVM & $r$ & Chặn trên chuẩn đặc trưng, với $K(\bm{x},\bm{x})\leq r^2$. \\
\hline
Boosting & $h_t$ & Bộ phân loại cơ sở/yếu tại vòng lặp $t$. \\
\hline
Boosting & $\alpha_t$ & Trọng số của bộ phân loại yếu $h_t$. \\
\hline
\end{tabular}
\end{table}

\begin{table}[H]
\centering
\small
\caption{Tổng hợp các ký hiệu chính trong báo cáo (tiếp)}
\begin{tabular}{|>{\raggedright\arraybackslash}p{3.2cm}|>{\raggedright\arraybackslash}p{4.2cm}|>{\raggedright\arraybackslash}p{6.6cm}|}
\hline
\textbf{Nhóm} & \textbf{Ký hiệu} & \textbf{Ý nghĩa} \\
\hline
Boosting & $g_n$ & Tổ hợp các bộ phân loại yếu sau $n$ vòng. \\
\hline
Boosting & $F(\bm{\alpha})$ & Hàm mục tiêu mũ cần tối ưu trong AdaBoost.MH. \\
\hline
Boosting & $D_t(i,l)$ & Trọng số của cặp mẫu--nhãn $(i,l)$ tại vòng $t$. \\
\hline
Cây quyết định & $n$ & Một nút trong cây quyết định. \\
\hline
Cây quyết định & $q$ & Câu hỏi phân tách tại nút, ví dụ $X_j\leq a$. \\
\hline
Cây quyết định & $p_l(n)$ & Tỷ lệ mẫu lớp $l$ rơi vào nút $n$. \\
\hline
Cây quyết định & $F(n)$ & Hàm đo độ vẩn đục tại nút $n$. \\
\hline
Cây quyết định & $\widetilde{F}(n,q)$ & Mức giảm độ vẩn đục khi tách nút $n$ bằng câu hỏi $q$. \\
\hline
Cây quyết định & $\lambda$ & Tham số điều chuẩn dùng để cắt tỉa cây. \\
\hline
ECOC/XMC & $M$ & Ma trận mã ECOC, mỗi hàng là codeword của một lớp. \\
\hline
\end{tabular}
\end{table}

\begin{table}[H]
\centering
\small
\caption{Tổng hợp các ký hiệu chính trong báo cáo (tiếp)}
\begin{tabular}{|>{\raggedright\arraybackslash}p{3.2cm}|>{\raggedright\arraybackslash}p{4.2cm}|>{\raggedright\arraybackslash}p{6.6cm}|}
\hline
\textbf{Nhóm} & \textbf{Ký hiệu} & \textbf{Ý nghĩa} \\
\hline
ECOC/XMC & $M_l$ & Codeword ứng với lớp $l$. \\
\hline
ECOC/XMC & $d_H$ & Khoảng cách Hamming giữa hai vector mã. \\
\hline
ECOC/XMC & $c$ & Độ dài mã ECOC, tương ứng số bộ phân loại nhị phân. \\
\hline
\end{tabular}
\end{table}

\subsection{Bảng tổng hợp công thức}

\renewcommand{\arraystretch}{1.25}
\setlength{\tabcolsep}{5pt}

\begin{table}[H]
\centering
\small
\caption{Tổng hợp các công thức chính trong báo cáo}
\begin{tabular}{|>{\raggedright\arraybackslash}p{3.2cm}|>{\raggedright\arraybackslash}p{5.6cm}|>{\raggedright\arraybackslash}p{5.2cm}|}
\hline
\textbf{Chủ đề} & \textbf{Công thức} & \textbf{Ý nghĩa} \\
\hline
Hàm điểm số
& $h:\mathcal{X}\times\mathcal{Y}\to\mathbb{R}$
& Gán điểm tin cậy cho cặp dữ liệu--nhãn $(\bm{x},y)$. \\
\hline
Quy tắc dự đoán đa lớp
& $g_h(\bm{x})=\arg\max_{y\in\mathcal{Y}}h(\bm{x},y)$
& Chọn nhãn có điểm số cao nhất. \\
\hline
Lề đa lớp
& $\rho_h(\bm{x},y)=h(\bm{x},y)-\max_{y'\neq y}h(\bm{x},y')$
& Khoảng cách giữa điểm nhãn đúng và nhãn sai cạnh tranh mạnh nhất. \\
\hline
Điều kiện phân loại sai
& $g_h(\bm{x})\neq y\Longleftrightarrow \rho_h(\bm{x},y)\leq 0$
& Mô hình sai khi lề không dương. \\
\hline
Sai số 0--1 kinh nghiệm
& $\widehat{R}_{0/1}(h)=\dfrac{1}{m}\sum_{i=1}^{m}\mathbb{1}[\rho_h(\bm{x}_i,y_i)\leq0]$
& Tỷ lệ mẫu huấn luyện bị phân loại sai. \\
\hline
Hàm mất mát lề
& $\Phi_\rho(t)=\begin{cases}1,&t\leq0\\1-t/\rho,&0<t<\rho\\0,&t\geq\rho\end{cases}$
& Phạt mẫu sai hoặc đúng nhưng chưa đạt lề mong muốn. \\
\hline
\end{tabular}
\end{table}

\begin{table}[H]
\centering
\small
\caption{Tổng hợp các công thức chính trong báo cáo (tiếp)}
\begin{tabular}{|>{\raggedright\arraybackslash}p{3.2cm}|>{\raggedright\arraybackslash}p{5.6cm}|>{\raggedright\arraybackslash}p{5.2cm}|}
\hline
\textbf{Chủ đề} & \textbf{Công thức} & \textbf{Ý nghĩa} \\
\hline
Sai số lề kinh nghiệm
& $\widehat{R}_\rho(h)=\dfrac{1}{m}\sum_{i=1}^{m}\Phi_\rho(\rho_h(\bm{x}_i,y_i))$
& Chặn trên mềm của sai số huấn luyện 0--1. \\
\hline
Độ phức tạp Rademacher
& $\widehat{\mathfrak{R}}_S(\mathcal{F})=\dfrac{1}{m}\mathbb{E}_{\bm{\sigma}}\left[\sup_{f\in\mathcal{F}}\sum_{i=1}^{m}\sigma_i f(\bm{x}_i)\right]$
& Đo khả năng lớp hàm khớp nhiễu ngẫu nhiên. \\
\hline
Bổ đề max
& $\widehat{\mathfrak{R}}_S(\mathcal{G})\leq\sum_{j=1}^{\ell}\widehat{\mathfrak{R}}_S(\mathcal{F}_j)$
& Chặn độ phức tạp của lớp hàm lấy cực đại. \\
\hline
Giới hạn tổng quát hóa đa lớp
& $R(h)\leq\widehat{R}_\rho(h)+\dfrac{2k^2}{\rho}\mathfrak{R}_m(\Pi_1(\mathcal{H}))+\sqrt{\dfrac{\log(1/\delta)}{2m}}$
& Chặn sai số thật bằng sai số lề, độ phức tạp mô hình và xác suất tin cậy. \\
\hline
Kernel bound
& $\mathfrak{R}_m(\Pi_1(\mathcal{H}_{K,p}))\leq\sqrt{\dfrac{r^2\Lambda^2}{m}}$
& Chặn Rademacher cho giả thuyết kernel khi $\|W\|\leq\Lambda$ và $K(\bm{x},\bm{x})\leq r^2$. \\
\hline
Hệ quả kernel
& $R(h)\leq\widehat{R}_\rho(h)+2k^2\sqrt{\dfrac{r^2\Lambda^2}{m\rho^2}}+\sqrt{\dfrac{\log(1/\delta)}{2m}}$
& Dạng cụ thể của giới hạn lề cho mô hình kernel. \\
\hline
\end{tabular}
\end{table}

\begin{table}[H]
\centering
\small
\caption{Tổng hợp các công thức chính trong báo cáo (tiếp)}
\begin{tabular}{|>{\raggedright\arraybackslash}p{3.2cm}|>{\raggedright\arraybackslash}p{5.6cm}|>{\raggedright\arraybackslash}p{5.2cm}|}
\hline
\textbf{Chủ đề} & \textbf{Công thức} & \textbf{Ý nghĩa} \\
\hline
Multi-class SVM primal
& $\min_{W,\bm{\xi}}\dfrac{1}{2}\sum_{l=1}^{k}\|\bm{w}_l\|^2+C\sum_{i=1}^{m}\xi_i$
& Tối ưu đồng thời lề lớn và lỗi huấn luyện nhỏ. \\
\hline
Ràng buộc SVM
& $\langle\bm{w}_{y_i},\Phi(\bm{x}_i)\rangle\geq\langle\bm{w}_{l},\Phi(\bm{x}_i)\rangle+1-\xi_i$
& Điểm nhãn đúng phải lớn hơn mọi nhãn sai ít nhất $1-\xi_i$. \\
\hline
OVA
& $h_{\text{OVA}}(\bm{x})=\arg\max_{l\in\mathcal{Y}}f_l(\bm{x})$
& Huấn luyện một bộ phân loại cho mỗi lớp so với phần còn lại. \\
\hline
Số mô hình OVO
& $\binom{k}{2}=\dfrac{k(k-1)}{2}$
& Số bộ phân loại cặp lớp cần huấn luyện. \\
\hline
ECOC decoding
& $h_{\text{ECOC}}(\bm{x})=\arg\min_{l\in\mathcal{Y}}d_H(M_l,\bm{h}(\bm{x}))$
& Chọn lớp có mã gần nhất với vector dự đoán. \\
\hline
Khoảng cách Hamming
& $d_H(\bm{u},\bm{v})=\dfrac{1}{2}\sum_{j=1}^{c}(1-u_jv_j)$
& Số bit khác nhau giữa hai codeword nhị phân. \\
\hline
\end{tabular}
\end{table}

\begin{table}[H]
\centering
\small
\caption{Tổng hợp các công thức chính trong báo cáo (tiếp)}
\begin{tabular}{|>{\raggedright\arraybackslash}p{3.2cm}|>{\raggedright\arraybackslash}p{5.6cm}|>{\raggedright\arraybackslash}p{5.2cm}|}
\hline
\textbf{Chủ đề} & \textbf{Công thức} & \textbf{Ý nghĩa} \\
\hline
AdaBoost.MH
& $F(\bm{\alpha})=\sum_{i=1}^{m}\sum_{l=1}^{k}\exp(-y_i[l]g_n(\bm{x}_i,l))$
& Hàm mục tiêu mũ cho bài toán đa nhãn. \\
\hline
Cập nhật trọng số Boosting
& $D_{t+1}(i,l)\propto D_t(i,l)\exp(-\alpha_t y_i[l]h_t(\bm{x}_i,l))$
& Tăng trọng số cho các cặp mẫu--nhãn dự đoán sai. \\
\hline
Entropy
& $\text{Entropy}(n)=-\sum_{l=1}^{k}p_l(n)\log p_l(n)$
& Độ vẩn đục của nút trong cây quyết định. \\
\hline
Gini index
& $\text{Gini}(n)=1-\sum_{l=1}^{k}p_l(n)^2$
& Độ vẩn đục phổ biến khi xây cây quyết định. \\
\hline
Cắt tỉa cây
& $G_\lambda(\text{tree})=\sum_{n\in\widetilde{\text{tree}}}|n|F(n)+\lambda|\widetilde{\text{tree}}|$
& Cân bằng giữa lỗi tại lá và độ phức tạp của cây. \\
\hline
\end{tabular}
\end{table}